
\begin{abstract}

无线电干扰源的快速自动定位与清除是无线电频谱管理的核心任务。本文针对机器狗搭载测向机在半径 1800 米的圆形目标区域内自主搜索并清除干扰源的问题，建立了一套从几何度量到序贯决策的完整数学模型，逐层解决了交会定位区域计算、第二检测点布站、全向源搜索清除以及含定向源的检测清除四个问题。

\textbf{对于问题一,}本文将示向度的 $\pm1^\circ$ 误差建模为以检测点为顶点、张角 $2^\circ$ 的凸楔形，并用两个半平面的交等价表示以规避方位角在 $0^\circ/360^\circ$ 处的回绕；定位区域即为全部楔形与目标圆域的交，采用 Sutherland--Hodgman 逐次半平面裁剪求其顶点，再用旋转卡壳在 $O(h)$ 时间内求凸多边形直径，并以顶点对暴力枚举互为校验，二者偏差不超过 $10^{-9}$ 米。针对覆盖性，本文由 Thales 逆定理给出充要判据：凸四边形被其直径圆覆盖当且仅当两个非直径端点处的内角均不小于 $90^\circ$。5000 组随机构型的统计表明直径圆覆盖率为 $98.99\%$，失效仅发生于交会角接近 $90^\circ$ 的临界构型（此时定位区域趋近矩形），且失败的绝对幅度极小，超出量平均仅 0.4928 米、最大 1.1697 米。

\textbf{对于问题二,}本文把单次示向度给出的可行域刻画为张角 $2^\circ$ 的楔形，推导了两条方位线在角扰动下的交点位移公式，并在垂直布站这一特例下得到完整闭式解：最优基线 $b^{*}=\sqrt2r$、最优交会角 $\alpha^{*}=\arccos(1/\sqrt3)=54.74^\circ$、最小定位误差 $E_{\min}=3\delta r=0.052360r$。针对目标距离未知的实际情形，本文在可行域上以期望定位误差为目标做二维寻优，得到最优布站为基线 1401.4 米、偏离示向度 $20.03^\circ$，期望定位误差 22.19 米；而教科书中常被引用的垂直布站即使把基线放长至 1510 米，误差仍高达 55.73 米，是最优方案的 2.5 倍。候选区域由目标函数的次水平集给出，为半径 $1230\sim1630$ 米、偏离示向度 $14^\circ\sim29^\circ$ 的两条对称圆弧带。

\textbf{对于问题三,}本文提出"全域粗探—交会定距—迭代逼近—就地清除"四阶段策略。粗探阶段的核心是保证覆盖网：由 Kershner 定理，6 个半径 1000 米的圆盘最多覆盖半径 1732.05 米的圆域，小于目标区域半径，故观测点数下界为 7；本文取"1 中心加 6 环点"构型，求得最小行程的环半径为 1122.96 米（覆盖半径恰为 1000 米），经 $1.5\times10^6$ 点稠密采样验证覆盖率 $99.999867\%$。由问题二结论反解得进入清除精度所需的逼近距离阈值为 381.97 米，据此设计迭代逼近流程。虚拟时间模型分解为移动、检测、切换与清除四项，并给出基于 Beardwood--Halton--Hammersley 定理的旅行商下界；$N=13$ 时模型总时间为 4450.3 秒、平均定位清除时间 342.3 秒，与理想下界 1847.0 秒的比值为 2.4。

\textbf{对于问题四,}定向源使可检测域由圆盘收缩为半圆盘，"无信号"退化为三义信息。本文证明了背向不可测定理（视角差不少于 $180^\circ$ 的两点不可能同时无信号），并给出一般性的方位包围判据：零漏测当且仅当观测点在干扰源有效接收半径内的方位角最大间隔小于 $180^\circ$。据此检验，问题三的七点覆盖网在半径 1200 米以外对定向源必然漏测，全域平均漏测概率 $65.24\%$；在"中心加分级同心环"构型族内求解得零漏测所需的最小规模为 37 点，其粗探时间为 8776.9 秒，相对问题三的惩罚系数为 3.025。为压低代价，本文进一步给出自适应补测策略：仅对未解频道沿加密环补测，在 2 个未解频道的典型情形下总时间 6561.8 秒，较静态方案节省 $25.2\%$。

最后，本文对模型做了交叉验证与灵敏度分析，讨论了其优点与不足，并给出了向多站测向组网、传感器盲区规避等场景推广的途径。

\keywords{交会定位；半平面求交；旋转卡壳；Thales 判据；Kershner 覆盖定理；旅行商下界；方位包围}

\end{abstract}
